\phantomsection\addcontentsline{toc}{section}{Future Transport}
\ECchapter{15}{Future Transport}{How will people and goods move tomorrow?}{ECOrange}

\ECObjectives{
\begin{itemize}
\item Compare transport modes using energy, capacity, space and infrastructure criteria.
\item Explain the perception--decision--action chain of an autonomous vehicle.
\item Design and justify a multimodal transport solution rather than a single vehicle technology.
\end{itemize}}

\begin{multicols}{2}
\begin{engineeringnews}
Battery buses, hydrogen trains, autonomous shuttles and electric aircraft are being tested around the
world. The largest environmental gains, however, may come from combining efficient vehicles with
public transport, active mobility, higher occupancy and better use of existing infrastructure.
\end{engineeringnews}

\begin{ecarticle}
Transport engineering must move people and goods safely, reliably and with limited energy and land
use. Because journey length, passenger capacity and infrastructure differ, no propulsion technology
is ideal for every application.

Battery-electric vehicles are efficient because electric motors convert a large proportion of stored
energy into motion and can recover energy during braking. They are well suited to urban buses, cars
and delivery vehicles that return regularly to a charging point. Long-distance freight requires
larger batteries, faster charging or another energy carrier. Hydrogen fuel cells may reduce onboard
storage mass in some heavy applications, but producing, compressing and converting hydrogen
introduces additional energy losses.

Rail carries many passengers or large freight loads with low rolling resistance. Electrified trains
can use energy directly from the grid and do not need to carry a large energy store. High-speed rail
can replace some short-haul flights when stations connect city centres and total door-to-door time is
competitive.

Autonomous vehicles add a digital engineering challenge. Cameras provide detailed images, radar
measures distance and relative speed, and lidar creates a three-dimensional representation. Sensor
fusion combines these signals. Software identifies objects, predicts their motion, plans a safe
trajectory and commands steering and braking. Adverse weather, unusual road layouts and human
behaviour make validation difficult.

Future mobility is therefore a system-level problem. A journey may combine walking, bicycle, train
and an on-demand shuttle. Ticketing, timetables and physical interchanges must work together.
Engineers must consider accessibility, affordability, charging networks, manufacture and the carbon
intensity of the energy supply. Replacing every conventional car with another type of car would not,
by itself, eliminate congestion, parking demand or unequal access to mobility.
\end{ecarticle}

\begin{tcolorbox}[title=\sffamily\bfseries Engineering Insight,colback=yellow!10,colframe=ECOrange]
Vehicle efficiency and transport efficiency are not identical. Occupancy, route, waiting time,
vehicle mass and infrastructure use can matter as much as propulsion efficiency. The correct unit of
comparison is often energy or emissions per passenger-kilometre or tonne-kilometre.
\end{tcolorbox}

\begin{engineeringfocus}
\textbf{Autonomous driving control chain}
\begin{center}
\resizebox{\linewidth}{!}{%
\begin{tikzpicture}[font=\scriptsize,>=Latex,node distance=3.5mm]
\node[draw=ECBlue,fill=ECLightBlue,rounded corners,align=center] (sen) {camera, radar,\\lidar, GNSS};
\node[draw=ECPurple,fill=ECPurple!10,rounded corners,align=center,right=of sen] (per) {sensor fusion\\and perception};
\node[draw=ECOrange,fill=ECOrange!10,rounded corners,align=center,right=of per] (plan) {prediction and\\trajectory planning};
\node[draw=ECGreen,fill=ECGreen!10,rounded corners,align=center,below=of plan] (act) {steering, motor\\and brakes};
\node[draw=ECRed,fill=ECRed!8,rounded corners,align=center,below=of per] (veh) {vehicle and\\road environment};
\draw[->,thick] (sen)--(per); \draw[->,thick] (per)--(plan); \draw[->,thick] (plan)--(act);
\draw[->,thick] (act)--(veh); \draw[->,thick,dashed] (veh)--(sen);
\end{tikzpicture}%
}
\end{center}
The loop must continue safely when a sensor is degraded. Redundancy, plausibility checks, driver or
remote supervision and a minimum-risk manoeuvre are central design features.
\end{engineeringfocus}

\begin{keyfigures}
\begin{tabularx}{\linewidth}{@{}Xr@{}}
Electric motor efficiency & high\\
Rail rolling resistance & low\\
Autonomous control loop & continuous\\
Useful comparison metric & passenger-km\\
Main urban constraint & limited space
\end{tabularx}
\end{keyfigures}

\begin{tcolorbox}[title=\sffamily\bfseries Illustrative operational emissions,colback=white,colframe=ECOrange]
\begin{center}
\begin{tikzpicture}
\begin{axis}[xbar,width=.93\linewidth,height=50mm,xlabel={Relative emissions per passenger-km},xmin=0,xmax=135,symbolic y coords={Bicycle,Electric rail,Rail,Coach,Car},ytick=data,ticklabel style={font=\scriptsize},grid=major]
\addplot table[x=emissions,y=mode,col sep=comma]{data/EC15/transport_emissions.csv};
\end{axis}
\end{tikzpicture}
\end{center}
\footnotesize Illustrative comparison only. Actual emissions depend strongly on occupancy, vehicle,
energy source, route and lifecycle boundaries.
\end{tcolorbox}

\begin{tcolorbox}[title=\sffamily\bfseries Multimodal journey,colback=white,colframe=ECBlue]
\begin{center}
\resizebox{\linewidth}{!}{%
\begin{tikzpicture}[font=\scriptsize,>=Latex,node distance=4mm]
\node[draw,rounded corners,fill=ECGreen!10] (home) {home};
\node[draw,rounded corners,fill=ECGreen!10,right=of home] (bike) {walking / bicycle};
\node[draw,rounded corners,fill=ECBlue!10,right=of bike] (rail) {rail};
\node[draw,rounded corners,fill=ECOrange!10,right=of rail] (shut) {shared shuttle};
\node[draw,rounded corners,fill=ECPurple!10,right=of shut] (dest) {destination};
\draw[->,thick] (home)--(bike); \draw[->,thick] (bike)--(rail); \draw[->,thick] (rail)--(shut); \draw[->,thick] (shut)--(dest);
\end{tikzpicture}%
}
\end{center}
A successful system minimises waiting, transfer distance, uncertainty and barriers for people with
reduced mobility.
\end{tcolorbox}

\begin{vocabulary}
\begin{tabularx}{\linewidth}{@{}lX@{}}
freight & fret\\
short-haul flight & vol court-courrier\\
rolling resistance & résistance au roulement\\
sensor fusion & fusion de capteurs\\
trajectory planning & planification de trajectoire\\
minimum-risk manoeuvre & manœuvre de risque minimal\\
shared mobility & mobilité partagée\\
interchange & pôle d'échanges\\
occupancy rate & taux de remplissage\\
passenger-kilometre & voyageur-kilomètre
\end{tabularx}
\end{vocabulary}

\begin{questions}
\item Why is no propulsion technology suitable for every journey?
\item Give one advantage of electrified rail over battery vehicles.
\item Describe the autonomous perception--decision--action loop.
\item Why is validation difficult in adverse or unusual conditions?
\item Use the graph to compare the highest- and lowest-emission motorised modes.
\item Why is vehicle occupancy important when comparing transport modes?
\item Which features make a multimodal journey attractive and accessible?
\end{questions}

\begin{applicationexercise}
A class of 30 students travels 18 km to a science museum. Compare five cars carrying six students
each with one coach carrying all 30 students, using the illustrative values in the graph.
\begin{enumerate}
\item Calculate the total passenger-kilometres for the trip.
\item Estimate the operational emissions for each option in relative units.
\item By what factor does the coach reduce emissions compared with the cars?
\item Recalculate the car option for an average occupancy of three students per car.
\item Give two reasons why a real comparison should also include lifecycle and infrastructure effects.
\end{enumerate}
\end{applicationexercise}

\begin{engineeringchallenge}
Propose a low-carbon transport plan between a suburban school and a city centre. Define demand and
accessibility requirements; compare at least three modes using capacity, journey time, energy,
emissions, land use, reliability and infrastructure; include first- and last-kilometre solutions;
and identify one disruption scenario with a fallback plan. Present a weighted decision matrix and a
justified recommendation.
\end{engineeringchallenge}

\begin{speaking}
Should city centres restrict private cars? Defend a position in a four-minute technical briefing
using transport capacity, accessibility, energy, public space and social impact, and propose one
realistic alternative package.
\end{speaking}
\end{multicols}

\subsection*{Sources}
\ECsource{International Energy Agency -- Transport}{https://www.iea.org/energy-system/transport}
\ECsource{European Commission -- Mobility and Transport}{https://transport.ec.europa.eu/index_en}
